%-------------------------------------------
%%%%%%% LUKAS FUCHS CHEATSHEAT %%%%%%%%%%%%%
%-------------------------------------------

\documentclass[8pt,a4paper,landscape]{extarticle}
\usepackage[
  landscape, 
  top=0.2cm,    % margin top
  bottom=0.2cm, % margin bottom
  left=0.2cm,   % margin left
  right=0.2cm   % margin right
]{geometry}
\usepackage{multicol}
\usepackage{amsmath}
\usepackage{amssymb}
\usepackage{xcolor}
\usepackage{caption}
\usepackage{enumitem}
\usepackage{mathtools}
\usepackage{titlesec}
\usepackage{tabularx}
\usepackage{array}
\usepackage{varwidth} %Important for calculations with \linewidth
\usepackage{enumitem}
\usepackage{soul}
\usepackage{graphicx}
\input{glyphtounicode}


% Ensure that generate PDF is machine readable/ATS parsable
\pdfgentounicode=1





%-------------------------------------------
%%%%%%  SETTINGS  %%%%%%%%%%%%%%%%%%%%%%%%%%
%-------------------------------------------

%----------Generell-------------------------
\pagestyle{empty}
\setlength{\columnsep}{10pt}
\setlength{\columnseprule}{0.4pt}
\setlength{\parskip}{0pt}
\setlength{\parindent}{0pt}
\raggedcolumns
\everymath{\thinmuskip=1mu \medmuskip=0mu \thickmuskip=1mu}

%Customize the sections
\titleformat{\section}
  {\vspace{-10pt}\bfseries\color{MyBlue}\raggedright\scshape\normalsize}
  {}
  {0em}
  {}
  [{\color{black}\titlerule\vspace{-5pt}}]

%---------Personalize-----------------------
\renewcommand{\columnseprulecolor}{\color{lightgray}}
%New costum colors
\definecolor{MyBoxColor}{rgb}{1,0.9,0.9}
\definecolor{MyBlue}{rgb}{0.01,0.5,0.8}
\definecolor{MyRed}{rgb}{0.8,0.01,0.01}
\definecolor{MyGray}{rgb}{0.2,0.2,0.2}
\definecolor{MyPurple1}{rgb}{0.7,0.6,1}
\definecolor{MyPurple2}{rgb}{0.4,0.2,0.5}

\setul{0.5ex}{0.6pt}

%New columtype: 'Y' as 'X' aligned to the left
\newcolumntype{Y}{>{\raggedright\arraybackslash}X}

%---------Custom Commands-------------------
% 1. Boxed Formula
% Command: \fboxit{formel}
\newcommand{\fboxit}[1]{
  \begingroup
    \setlength{\fboxsep}{1.5pt}
    \colorbox{MyPurple1}{
      \begin{varwidth}{\dimexpr\linewidth-5\fboxsep\relax} %sets maximum width on coloum width
        \ensuremath{\displaystyle #1}
      \end{varwidth}
    }
  \endgroup
}

% 2. Compact List/Table
% Command: \fgridtable{Titel}{Spaltenanzahl}{Inhalt}
\newcommand{\fgridtable}[3]{
  \vspace{0pt}
    {
    \par
    \textbf{#1:} \\
    \renewcommand{\arraystretch}{1.3} % Increases the row hight for better readability
    \setlength{\tabcolsep}{2pt}      % Minimal margin to the side
    \begin{tabularx}{\linewidth}{|#2|} % \linewidth tells table to stay in his coloum
      \hline
      #3 \\ \hline
    \end{tabularx}
    }
  \vspace{-1pt}
}
%Example:
%\fgridtable{Vektoren}{Y|Y|Y}{
% x & y & z \\ \hline
% 1 & 0 & 0 \\ \hline
% 0 & 1 & 0
%}

% 3. Simple Alignment
%Command: \falignleft{formulas}
\newcommand{\falign}[1]{
  \vspace{2pt}
  {
    \noindent
    $\begin{aligned}
      #1
    \end{aligned}$
    \hspace{-2pt}
  }
  \vspace{0pt}
}
%Example: 
%\falign{
%    f(x) &= x^2 \\
%    g(x) &= \sin(x) \\
%    a + c &= \text{Ergebnis} \\
%    2x + 15 &= 35 \\
%    x &= 10
%}
\newcommand{\faligntext}[2]{%
  \par\noindent
  \vspace{3pt}
  \begin{tabularx}{\linewidth}{@{}l X@{}} % l = Formulawiidth, X = Free space
    \vspace{0pt}$\begin{aligned}[t] #1 \end{aligned}$ & 
    \vspace{0pt}\raggedright\tiny #2 
  \end{tabularx}
  \par\vspace{-2pt}
}

%4. Image inculde command
% Command: \fimagetext{Width percantage}{Image name}{Text}
\newcommand{\fimagetext}[3]{%
  \vspace{2pt}
  \begin{tabularx}{\linewidth}{@{} p{#1\linewidth} X @{}}
    \vspace{0pt}\includegraphics[width=\linewidth]{#2} & 
    \vspace{0pt}\raggedright\tiny #3 
  \end{tabularx}
  \par\vspace{-2pt}
}
%Example:
%\fimagetext{0.4}{mein_diagramm.png}{some text at the side}


% 5. Heading and Item commands
%i)
\newcommand{\CheatsheetItem}[3]{
  \vspace{2pt}
  \item{
    \small\tiny\textbf{#1}{#2 \vspace{7pt}}{#3}
  }
}

%ii)
\newcommand{\CheatsheetHeading}[2]{
  \vspace{1pt}
    \begin{tabular*}{0.99\textwidth}[t]{@{}l@{\extracolsep{\fill}}r}
      {\setulcolor{MyPurple2}\color{MyPurple2}\ul{\textbf{#1}}} \textit{\color{gray}#2} \\
    \end{tabular*}\vspace{2pt}
}

\newcommand{\CheatsheetSubHeading}[2]{
  \vspace{1pt}
    \begin{tabular*}{0.99\textwidth}[t]{@{}l@{\extracolsep{\fill}}r}
      {\setulcolor{orange}\color{orange}\ul{\tiny\textbf{#1}}} \textit{\color{gray}#2} \\
    \end{tabular*}\vspace{1pt}
}

\renewcommand{\labelitemi}{\textbullet} 
%\renewcommand{\labelitemi}{$\triangleright$}
%\renewcommand{\labelitemi}{} 

\newcommand{\CheatsheetItemListStart}{\begin{itemize}[leftmargin=*, labelsep=5pt, nosep]}
\newcommand{\CheatsheetItemListEnd}{\end{itemize}\vspace{-5pt}}
% End of Heading and Items






%-------------------------------------------
%%%%%%  CHEATSHEET STARTS HERE  %%%%%%%%%%%%
%-------------------------------------------
\begin{document}
\tiny\color{MyGray}
\begin{multicols*}{5}

%---------First Row-------------------------
{\huge\color{MyPurple1}\textbf{Stochastische Signale}}
\vspace{4pt}
\section*{Wahrscheinlichkeitsräume:}
Wahrscheinlichkeitsraum $(\Omega, \mathbb{F}, P)$ besteht aus:
  \vspace{1pt}
  \CheatsheetHeading{Ergebnismenge $\Omega$: }{}
  $\Omega = \{\omega_1, \omega_2, \dots\}$: \text{Menge aller möglichen Ergebnisse.}
  
  \vspace{1pt}
  \CheatsheetHeading{Ereignisalgebra $\mathbb{F}$:}{}
  $\mathbb{F} = \{A_1, A_2, \dots\}$: \text{Menge von Ereignissen} $A_i \subseteq \Omega$.

    \vspace{0pt}
    \CheatsheetItemListStart
      \CheatsheetItem{Eigneschaften: }{}{}
    \CheatsheetItemListEnd

      \vspace{0pt}
      $\vspace{3pt}\text{- } \Omega \in \mathbb{F} $ \\
      $\text{- } A_i \in \mathbb{F} \ \Rightarrow \ A_i^c \in \mathbb{F} $ \\
      $\text{- } A_1, \dots, A_k \in \mathbb{F} \ \Rightarrow \ \bigcup_{i=1}^{k} A_i \in \mathbb{F} $ \\
      %
      \text{- Impliziert: }$\text{- }\emptyset \in \mathbb{F} \text{ }(\Omega^c = \emptyset)\text{ }$ 
      $\text{- } A_i\backslash A_k \in \mathbb{F}\text{ }$
      $\text{- } \bigcap_{i = 1}^k \in \mathbb{F}$
      %
    \vspace{0pt}
    \CheatsheetItemListStart
     \CheatsheetItem{Größe: }{}{}
    \CheatsheetItemListEnd 

      \vspace{0pt}
      $\text{- }|\mathbb{F}| \hat{=} 2^{\text{Anzahl an disjunkter Teilmengen}}$ 

  \vspace{1pt}
  \CheatsheetHeading{Wahrscheinlichkeitsmaß P:}{}
    \CheatsheetItemListStart
        \CheatsheetItem{Axiome von Kolmogorow: }{}{}
    \CheatsheetItemListEnd
    
      \vspace{0pt}
      $\vspace{3pt}\text{- } P(A) \ge 0 \ \Rightarrow \ P: \mathbb{F} \mapsto [0,1] \\$
      $\vspace{-8pt}\text{- }P(\Omega) = 1$ \\
      $\text{- }P\left(\bigcup A_i\right) = \sum P(A_i) \quad \text{wenn: } \overbrace{A_i \bigcap A_j = \emptyset, \forall i \neq j}^{\text{(disjunkt)}}$

  \vspace{1pt}
  \CheatsheetHeading{Grundlegende Rechenregeln:}{}
    \vspace{-3pt}
    
    $\vspace{3pt}\text{- }P(A^c) = 1 - P(A) $ \\
    $\vspace{3pt}\text{- }P(\emptyset) = 0 $ \\
    $\vspace{3pt}\text{- }P(A \setminus B) = P(A) - P(A \cap B) $ \\
    $\vspace{3pt}\text{- }P(A \cup B) = P(A) + P(B) - P(A \cap B)$
    

    



%---------Second Row------------------------
\columnbreak

%---------Third Row-------------------------
\columnbreak


%---------Fourth Row------------------------
\columnbreak


%---------Fifth Row-------------------------
\columnbreak


%---------Sixth Row-------------------------
\columnbreak


%-------------------------------------------
%%%%%%  CHEATSHEET ENDS HERE  %%%%%%%%%%%%
%-------------------------------------------

\end{multicols*}
\end{document}